%%%%%%%%%%%%%%%%%%%%%%%%%%%%%%%%%%%%%%%%%%%%%%%%%%%%%%%%%%%%%%%%%%%%%%%%%%%%%
\subsection{Pipeline Hazard Handling (\texttt{RV\_PipelineCacheFwd} only)}
\label{subsec:pipehazard}
%%%%%%%%%%%%%%%%%%%%%%%%%%%%%%%%%%%%%%%%%%%%%%%%%%%%%%%%%%%%%%%%%%%%%%%%%%%%%

\texttt{RV\_NoPipeline} has no hazards to resolve (every instruction
completes in a single cycle before the next one starts); this section
applies only to \texttt{RV\_PipelineCacheFwd}'s 5-stage pipeline
(\texttt{cc\_cpupipefwd.sv}, IF/ID/EX/MEM/WB).

\subsubsection{Forwarding Paths}

Two forwarding paths resolve read-after-write (RAW) hazards without
stalling, feeding the EX-stage ALU operand muxes directly from later
pipeline stages instead of the (stale) value latched in the ID/EX
register:

\begin{itemize}
  \item \textbf{EX/MEM $\rightarrow$ EX} (1 instruction back): the
    ALU result (or PC+4 / CSR read value) held in \texttt{ex\_mem\_aluRes\_r}
    / \texttt{ex\_mem\_csr\_rdata\_r} / \texttt{ex\_mem\_pc\_r} is forwarded
    when the MEM-stage instruction writes a register the EX-stage
    instruction reads. Excluded for loads: load data is not yet available
    at this point (see load-use stall below).
  \item \textbf{MEM/WB $\rightarrow$ EX} (2 instructions back): the final
    write-back value \texttt{wb\_result\_s} (already muxed between ALU
    result, load data, PC+4, and CSR read value) is forwarded when the
    WB-stage instruction writes a register the EX-stage instruction reads.
    This path also carries load data, so a load-to-use hazard 2 instructions
    apart is resolved without a stall.
\end{itemize}

A third RAW distance -- \textbf{k=3} (producer in WB, consumer in ID, same
clock edge) -- needs no explicit forwarding mux: \texttt{as\_regfile} is
instantiated with \texttt{WBR=1}, which bypasses the write data
combinatorially to the read ports within the same cycle, so \texttt{id\_regA\_s}/
\texttt{id\_regB\_s} already see the new value when both events coincide.

\subsubsection{Remaining Stall Sources}

Two hazard classes cannot be resolved by forwarding and still stall the
pipeline (freeze all stages via \texttt{global\_stall\_s}, or insert a
bubble via \texttt{data\_stall\_s}):

\begin{itemize}
  \item \textbf{Load-use (1 cycle):} a \texttt{LOAD} is in EX
    (\texttt{id\_ex\_dMemRd\_r}=1) and the very next instruction (in ID)
    reads the register it will load into. The load's result is only ready
    at the \emph{end} of the following MEM stage, too late to forward to a
    consumer already in EX at the same time. The pipeline inserts one
    bubble (NOP into EX, consumer held in ID); the next cycle the load
    reaches WB and the consumer reaches EX, where MEM/WB $\rightarrow$ EX
    forwarding resolves it.
  \item \textbf{CSR hazard (conservative):} CSR writes commit at WB and
    there is no forwarding path for CSR register values. Any CSR
    instruction in ID is held (\texttt{data\_stall\_s}) until no earlier
    CSR instruction remains anywhere in the pipeline.
\end{itemize}

\subsubsection{Control Hazards}

A taken branch or an unconditional jump (JAL/JALR) is only resolved in EX,
after the two following instructions have already been speculatively
fetched into IF and ID. \texttt{ex\_branch\_taken\_s} (branch) and
\texttt{wb\_trap\_taken\_s}/\texttt{wb\_mret\_s} (trap/\texttt{mret})
each assert \texttt{flush\_s}, which squashes the IF/ID and ID/EX pipeline
registers (2-cycle bubble) regardless of \texttt{global\_stall\_s} -- see
Section~\ref{subsubsec:flushStallRace} for why that last part specifically
had to be fixed.

\subsubsection{Hazard Test Suite}
\label{subsubsec:hazardTestSuite}

Five directed ASM tests (\texttt{rv64Sim/hazard\_*.asm}, run under
\texttt{tb\_rv64i\_instr.sv} like the base ISA tests) exercise every
hazard class above:

\begin{table}[H]
\caption{Pipeline Hazard Test Suite}
\label{tab:hazardTests}
\centering
\begin{tabularx}{\textwidth}{|l |X|}
  \hline
  Test & Coverage \\
  \hline
  \hline
  \texttt{hazard\_raw\_alu} & ALU$\to$ALU RAW at distance k=1
    (EX/MEM$\to$EX), k=2 (MEM/WB$\to$EX), k=3 (regfile WBR bypass);
    no stall expected for any of the three. \\
  \hline
  \texttt{hazard\_raw\_load} & Three load-use RAW hazards; each forces
    exactly one 1-cycle stall. \\
  \hline
  \texttt{hazard\_waw\_war} & Write-after-write (second write wins) and
    write-after-read (read observes the old value); both are naturally
    transparent in this in-order, single-issue pipeline. \\
  \hline
  \texttt{hazard\_ctrl\_branch} & Two taken branches (2-cycle flush each)
    plus one not-taken branch (no flush). \\
  \hline
  \texttt{hazard\_ctrl\_jump} & Forward \texttt{JAL} flush, backward
    \texttt{JAL} loop, and a \texttt{JAL}+\texttt{JALR} subroutine
    call/return pair. \\
  \hline
\end{tabularx}
\end{table}

A separate statistics testbench, \texttt{tb\_rv64i\_hazard\_stat.sv},
re-runs each of the five programs and counts \texttt{clk\_div\_s} cycles
from reset de-assertion to the program's final (successful) GPIO write,
reporting \texttt{cycles}, \texttt{instr}, and \texttt{CPI} (see the
regression log excerpt in Section~\ref{sec:simtest01}); this is a
performance/regression indicator, not a pass/fail correctness check --
correctness is still the GPIO\,=\,0x55 comparison every test in this suite
uses.
